\noindent This Supplementary Information provides implementation details and audit tables supporting the manuscript. The main benchmark results use unconstrained pocket-conditioned generation. An optional mode generates the unfixed portion of a molecule while retaining specified atom positions and element types.

\section*{Supplementary Note 1 | Generation with fixed atom subsets}

NA2smol includes an optional fixed-subset generation mode that retains specified atom positions and types (Algorithm~\ref{alg:constrained}). Reported benchmarks and the KRAS G-quadruplex case study use generation from the target pocket without fixed ligand fragments (Methods).

For scaffold-based redesign, this mode conditions generation on a subset of atoms with fixed positions and types (the fixed subset). The fixed atoms are made to follow the same forward Bayesian-flow dynamics used during training, whereas the remaining (free) atoms are generated through the standard reverse sampling process. Algorithm~\ref{alg:constrained} gives the corresponding fixed-code procedure.

\begin{algorithm}[H]
\caption{Seed-constrained inpainting sampling}
\label{alg:constrained}
\begin{algorithmic}[1]
\Require Fixed atom mask $\mathbf{m} \in \{0,1\}^N$, fixed positions $\mathbf{X}_{\text{fixed}}$, fixed types $\mathbf{T}_{\text{fixed}}$ (one-hot), sampling steps $T$, pocket conditioning $\mathcal{P}$
\Ensure Generated ligand whose seed atoms equal $(\mathbf{X}_{\text{fixed}}, \mathbf{T}_{\text{fixed}})$

\Statex \textbf{Phase 1: Pre-compute forward states for fixed atoms}

\For{$i \gets 0$ to $T$}
    \State $t_i \gets i / T$
    \State $\gamma_i \gets 1 - \sigma_1^{2t_i}$
    \State $\boldsymbol{\mu}_{\text{coord}}^{\text{fixed}}[i, \mathbf{m}] \gets \gamma_i \cdot \mathbf{X}_{\text{fixed}}$
    \State $\boldsymbol{\theta}_{\text{type}}^{\text{fixed}}[i, \mathbf{m}] \gets \mathbf{T}_{\text{fixed}}$
\EndFor

\Statex \textbf{Phase 2: Initialise free atoms}

\State $\boldsymbol{\mu}_{\text{coord}}^{(0)}[\neg \mathbf{m}] \gets \mathbf{0}$
\State $\boldsymbol{\theta}_{\text{type}}^{(0)}[\neg \mathbf{m}] \gets \mathbf{1}/K$

\Statex \textbf{Phase 3: Joint reverse sampling with fixed-atom states}

\For{$i \gets 1$ to $T$}
    \State $t \gets (i-1) / T$,\quad $t^{+} \gets i / T$

    \State \textbf{Assemble input coordinate and atom-type distributions} (fixed atoms from the precomputed forward states, free atoms from the running state):
    \State $\boldsymbol{\mu}_{\text{coord}}^{\text{input}} \gets \boldsymbol{\mu}_{\text{coord}}^{\text{fixed}}[i-1, \mathbf{m}]$ on $\mathbf{m}$; $\boldsymbol{\mu}_{\text{coord}}^{(i-1)}[\neg \mathbf{m}]$ elsewhere
    \State $\boldsymbol{\theta}_{\text{type}}^{\text{input}} \gets \boldsymbol{\theta}_{\text{type}}^{\text{fixed}}[i-1, \mathbf{m}]$ on $\mathbf{m}$; $\boldsymbol{\theta}_{\text{type}}^{(i-1)}[\neg \mathbf{m}]$ elsewhere

    \State \textbf{Model prediction on the full graph} (fixed atoms included):
    \State $\hat{\mathbf{X}}^{(i)}, \hat{\mathbf{p}}_0^{(i)} \gets \text{NA2smol}(\boldsymbol{\mu}_{\text{coord}}^{\text{input}}, \boldsymbol{\theta}_{\text{type}}^{\text{input}}, t, \mathcal{P})$

    \State \textbf{Sample atom types} for output bookkeeping:
    \State $k \sim \mathrm{Categorical}(\hat{\mathbf{p}}_0^{(i)})$;\quad $k[\mathbf{m}] \gets$ fixed classes of $\mathbf{T}_{\text{fixed}}$

    \State \textbf{Closed-form Bayesian-flow update of free atoms at $t^{+}$:}
    \State $\gamma^{+} \gets 1 - \sigma_1^{2t^{+}}$
    \State $\boldsymbol{\mu}_{\text{coord}}^{(i)}[\neg \mathbf{m}] \gets \gamma^{+} \hat{\mathbf{X}}^{(i)}[\neg \mathbf{m}] + \sqrt{\gamma^{+}(1-\gamma^{+})}\, \boldsymbol{\epsilon}_{\text{coord}},\ \ \boldsymbol{\epsilon}_{\text{coord}} \sim \mathcal{N}(\mathbf{0}, \mathbf{I})$
    \State $\boldsymbol{\theta}_{\text{type}}^{(i)}[\neg \mathbf{m}] \gets \mathrm{softmax}\bigl(\beta_1 (t^{+})^2 (K \hat{\mathbf{p}}_0^{(i)}[\neg \mathbf{m}] - \mathbf{1}) + \sqrt{\beta_1 (t^{+})^2 K}\, \boldsymbol{\epsilon}_{\text{type}}\bigr),\ \ \boldsymbol{\epsilon}_{\text{type}} \sim \mathcal{N}(\mathbf{0}, \mathbf{I})$

    \State \textbf{Fixed-subset update:} fixed atoms are excluded from both updates (their states enter only via the assembled inputs).
\EndFor

\State \textbf{Final prediction at $t = 1$} with inputs assembled from $\boldsymbol{\mu}^{\text{fixed}}[T, \mathbf{m}]$, $\boldsymbol{\theta}^{\text{fixed}}[T, \mathbf{m}]$ and the final free-atom states:
\State $\hat{\mathbf{X}}_{\text{final}}, \hat{\mathbf{p}}_{0,\text{final}} \gets \text{NA2smol}(\boldsymbol{\mu}_{\text{coord}}^{\text{input}}, \boldsymbol{\theta}_{\text{type}}^{\text{input}}, 1, \mathcal{P})$

\State \Return $\hat{\mathbf{X}}_{\text{final}}$, $\hat{\mathbf{p}}_{0,\text{final}}$ (seed atoms read out as $\mathbf{X}_{\text{fixed}}$, $\mathbf{T}_{\text{fixed}}$)
\end{algorithmic}
\end{algorithm}

Fixed atoms follow precomputed forward trajectories: coordinates follow $\boldsymbol{\mu}(t) = \gamma(t) \cdot \mathbf{X}_{\text{fixed}}$ with $\gamma(t) = 1 - \sigma_1^{2t}$, and atom-type distributions are held at the one-hot fixed-subset types. These states are inserted into the sampler at every step and overwrite any update to the fixed atoms; the sampled class of a fixed atom is likewise overwritten with its seed class. Free atoms are initialised from the prior ($\boldsymbol{\mu} = \mathbf{0}$, $\boldsymbol{\theta} = \mathbf{1}/K$) and are refreshed at each step by the closed-form Bayesian-flow updates applied to the model's own prediction $(\hat{\mathbf{X}}, \hat{\mathbf{p}}_0)$ at $t^{+}$. A variant that instead accumulates message-passing updates with precision $\rho_{\text{coord}}$ (the update style sketched in the main-text sampling algorithm) is also implemented; both variants share Phases 1, 2 and the fixed-atom-state structure above. The model processes the complete graph of fixed and free atoms at every step, allowing the fixed subset to condition the placement and atom types of the generated remainder.

\section*{Supplementary Note 2 | Reproducibility and implementation details}

Code and selected metadata are available in the repository listed in Data availability. Provenance records, reference tables and structure-identifier cross-checks support reproducibility. Large training databases, model checkpoints and benchmark outputs require separate deposition.

\section*{Supplementary Note 3 | Provenance of the fragment-based KRAS G-quadruplex ligand}

The 2,2$'$-methylenebis(benzothiazole) molecule was produced by the fragment-recombination workflow (Fig.~4); the crystallographic berberine ligand was not used as a template. Its benzothiazole ring systems derive from the pooled ring fragments and their heteroatom variants, joined through a matched attachment site (same removed-neighbour element; neighbour-position separation $\leq1.0$~\AA, Methods). The molecule's BER102 pose is the relaxed multi-seed docking pose from this workflow. Because the fragment pool was generated from the BER102 pocket, the assembled pose occupies the same site as the crystallographic ligand.

\section*{Supplementary Note 4 | Forward-pass specification of the evaluated configuration}

The teacher and the student have separate parameters; in the fine-tuning configuration both are initialised from checkpoint step 369999. Teacher calls use electron-density conditioning with empty pocket inputs, and student calls use pocket conditioning with empty electron-density inputs. Both receive the same noisy ligand state.

\begin{algorithm}[H]
\caption{Constrained training step as implemented}
\begin{algorithmic}[1]
\Require Ligand coordinates $X$, types $V$, pocket $P$, simulated ligand density $E$, timestep $t$
\State Generate one noisy ligand state $(\mu_t,\theta_t)$ using the Bayesian-flow updates
\State With gradients disabled, compute $(X_T,p_T,h_T)\gets f_T(\mu_t,\theta_t,t,E,\varnothing)$
\State Compute $(X_S,p_S,h_S)\gets f_S(\mu_t,\theta_t,t,\varnothing,P)$
\State \textbf{Output alignment.} $L_x\gets\operatorname{mean}_i\|X_{T,i}-X_{S,i}\|^2$ and $L_p\gets\frac12\operatorname{mean}_i[\mathrm{KL}(p_{T,i}\|p_{S,i})+\mathrm{KL}(p_{S,i}\|p_{T,i})]$
\State $w\gets\operatorname{sigmoid}(10(t-0.5))$, which increases with $t$ so the alignment acts most strongly as $t\to1$
\State \textbf{Contrastive term.} Pool $h_S$ and $h_T$ into molecule-level representations $z_S,z_T$ by attention pooling, project both through a shared projection head and $L_2$-normalise
\State $L_{\mathrm{contrastive}}\gets$ symmetric InfoNCE between $\{z_S\}$ and $\{z_T\}$, with the student and teacher views of the same molecule as the positive pair
\State $L\gets\operatorname{mean}_i[w_i(L_{\mathrm{coord},i}+L_{\mathrm{type},i})]+0.1\,w\,(L_x+L_p)+0.1\,L_{\mathrm{contrastive}}$
\State Update student parameters only; the teacher is frozen
\end{algorithmic}
\end{algorithm}

The teacher--student constraint therefore acts on the two paths' \emph{outputs}: the predicted coordinates and the predicted atom-type distributions. The contrastive term instead acts on pooled molecule-level representations of the two paths, and corresponds to the $d(z_{\mathrm{ED}},z_{\mathrm{lig}})$ distance drawn in the latent-space panel of Fig.~1. It is the objective that distinguishes the full-model configuration from Pretrain+Constrain, which applies the alignment term only.


\clearpage

\section*{Supplementary Note 5 | Chemical-overlap audit and recovery sensitivity}

We quantified chemical overlap between benchmark reference actives, the fine-tuning training partition and the ligand-only pretraining library to establish how much of the retrospective ECFP4-Tanimoto recovery could be explained by chemical familiarity rather than by pocket-conditioned generation. Per-molecule counts and percentages are in Supplementary Table~5, and recovery values under each exclusion are in Extended Data Tables~6 and~7; this note summarises what they show.

Identity between NA2smol outputs and fine-tuning training ligands is rare, so the fine-tuning split does not explain the recovery signal. Higher similarity to the pretraining library is expected, as pretraining exposes the model to a broad chemical space, and targets with reference actives in that space are the ones the model is designed to handle. Excluding targets with actives resembling the fine-tuning partition leaves NA2smol ahead on both benchmarks, confirming the fine-tuning split does not drive the ranking. Excluding targets with actives resembling the pretraining library changes the order: NA2smol falls to 0.053 on RNA (ED2Mol: 0.085) and to 0.238 on DUD-E (DiffSBDD: 0.455). DiffSBDD's DUD-E recovery remains nearly unaffected (0.447 to 0.455). These shifts follow from removing targets that the pretraining stage was designed to cover.

Excluding reference actives by structure identifier rather than by record yields the same results: the eligible-target set shrinks, but the recovery comparison between methods remains unchanged. These sensitivity analyses support the reported benchmark.

\section*{Supplementary Note 6 | Target specificity of ECFP4 recovery}
Recovery can be decomposed into the part explained by generic similarity to reference actives and
the part that depends on which target a molecule was generated for. We computed the difference
between the observed recovery and its value under shuffled molecule-to-target assignment (1,000
permutations), which isolates the second part. On the protein benchmark the largest difference
was for DiffSBDD (0.077 against a null of 0.380) followed by NA2smol (0.019 against 0.637); on the
RNA test set NA2smol gave the largest (0.037 against 0.180), with the remaining methods between
$-0.076$ and $0.019$. These differences are small for every method, and we do not read them as
evidence of a difference in target selectivity. Separating selectivity from the composition of the
generated library would require comparing the similarity of molecules generated for different
targets with that of the reference ligands of those targets, which this benchmark does not provide.
The quantity is reported here for completeness and is not used as an axis in Figs.~2
or 3.
